\section*{Chapitre \ref{ch:b1:lab-techniques-1} --- Techniques de laboratoire, I~: sécurité, mesure et séparation}

\begin{solution}{exo:b1:lab-techniques-1:1}
\emph{Danger}~: H225 (liquide très inflammable, catégorie 2) l'impose, et
l'étiquette porte la mention la plus grave. H225 est un \omterm{def:b1:lab-techniques-1:hazard}{danger}
physique~; H319 (sévère irritation des yeux) et H336 (somnolence ou
vertiges) sont des \omterm{def:b1:lab-techniques-1:hazard}{dangers} pour la santé. Précautions~: ni flamme ni
étincelle à proximité, un flacon fermé~; des lunettes~; travailler dans
un local ventilé ou sous la hotte.
% ledger: ghs:acetone, hstat:H225, hstat:H319, hstat:H336
\end{solution}

\begin{solution}{exo:b1:lab-techniques-1:2}
(a) Un \omterm{def:b1:lab-techniques-1:hazard}{danger}, propriété de l'acide. (b) Un \omterm{def:b1:lab-techniques-1:hazard}{risque}, faible parce que
l'exposition est faible. (c) Un \omterm{def:b1:lab-techniques-1:hazard}{danger}. (d) Un \omterm{def:b1:lab-techniques-1:hazard}{risque}~: même \omterm{def:b1:lab-techniques-1:hazard}{danger},
exposition plus grande (vapeurs, projections) à chaud et en récipient
ouvert.
\end{solution}

\begin{solution}{exo:b1:lab-techniques-1:3}
$R_f = 2.1/6.0 = 0.35$ et $4.2/6.0 = 0.70$. La seconde tache a le même
$R_f$ que la référence~: l'échantillon peut la contenir (c'est
cohérent), mais des $R_f$ égaux avec un seul éluant ne prouvent pas
l'identité~; un second éluant, ou un co-dépôt de l'échantillon et de la
référence qui reste une tache unique, renforce la conclusion. La tache à
0.35 est à coup sûr une autre espèce.
\end{solution}

\begin{solution}{exo:b1:lab-techniques-1:4}
Demi-largeur \qty{0.05}{mL}, distribution rectangulaire~: $u = 0.05/\sqrt3 = \qty{0.029}{mL}$.
Pour une différence de deux lectures indépendantes,
$u = \sqrt{2} \times 0.029 = \qty{0.041}{mL}$.
\end{solution}

\begin{solution}{exo:b1:lab-techniques-1:5}
Moyenne \qty{0.25100}{g}~; écarts $+2$, $-2$, $+5$, 0, $-5$ (en
\qty{e-4}{g}), somme des carrés \qty{58e-8}{g^2}, $s = \sqrt{58 \times
10^{-8}/4} = \qty{3.8e-4}{g}$~; $u = s/\sqrt5 = \qty{1.7e-4}{g}$~;
$U = \qty{3.4e-4}{g}$~: $m = (0.25100 \pm 0.00034)$~g, $k = 2$.
\end{solution}

\begin{solution}{exo:b1:lab-techniques-1:6}
Une portion~: $f = 1/(1 + 3 \times 30/50) = 0.357$, soit
\qty{64.3}{\percent} extraits. Deux portions de \qty{15}{mL}~:
$f = 1.9^{-2} = 0.277$,
\qty{72.3}{\percent}. Trois de \qty{10}{mL}~: $f = 1.6^{-3} = 0.244$,
\qty{75.6}{\percent}.
\end{solution}

\begin{solution}{exo:b1:lab-techniques-1:7}
$z = 0.0012/\sqrt{0.0006^2 + 0.0002^2} = 0.0012/0.00063 = 1.9 \le 2$~:
compatibles, de justesse. Avec \qty{0.0004}{mol/L}~: $z = 0.0012/0.00045 = 2.7$,
non compatibles~: le \omterm{def:b1:titration-methods:titration}{titrage} plus précis révèle une différence réelle
(une erreur de préparation, ou de \omterm{def:b1:titration-methods:titration}{titrage}).
\end{solution}

\begin{solution}{exo:b1:lab-techniques-1:8}
Q~: peu soluble à froid, très soluble à chaud. P ne dissout guère mieux
le solide à chaud qu'à froid (rien ne cristallise au refroidissement)~;
R en garde trop en solution à froid (pertes importantes). Avec Q,
\qty{5.0}{g} exigent au moins
$5.0/12 \times 100 = \qty{42}{mL}$ de solvant bouillant, qui retiennent
$0.3 \times 0.417 = \qty{0.13}{g}$ en solution à froid~: on peut
récupérer au plus
\qty{4.9}{g}, soit \qty{97.5}{\percent} (avant toute perte sur le
filtre).
\end{solution}

\begin{solution}{exo:b1:lab-techniques-1:9}
L'éthanol ($n_D = 1.3611$ à \qty{20}{\celsius})~; l'eau (1.333) et le
cyclohexane (1.4266) sont exclus. L'indice varie avec la température~:
une valeur ne signifie rien sans elle. Une valeur de 1.350, comprise
entre celles de l'eau et de l'éthanol, suggère un mélange, par exemple
de l'éthanol contenant de l'eau.
% ledger: nD:water, nD:ethanol, nD:cyclohexane
\end{solution}

\begin{solution}{exo:b1:lab-techniques-1:10}
Avec $t = KV/V_a$ et $n$ portions, $\ln f_n = -n\ln(1 + t/n)$~; quand
$n \to \infty$, $n \ln(1 + t/n) \to t$ (car $\ln(1 + \varepsilon) \sim
\varepsilon$), et $f_n$ décroît vers $\mathrm{e}^{-t}$ (\cref{prop:b1:lab-techniques-1:multiple-extractions}).
Ici $t = 3 \times 30/50 = 1.8$, $\mathrm{e}^{-1.8} = 0.165$~: au plus
\qty{83.5}{\percent}. Trois portions donnent déjà \qty{75.6}{\percent},
soit \qty{91}{\percent} de la limite~; en atteindre \qty{99}{\percent}
exige 32 portions de moins de \qty{1}{mL}~: au-delà de deux ou trois
portions, le gain ne vaut pas la peine.
\end{solution}

\begin{solution}{exo:b1:lab-techniques-1:11}
Avec $t = x - x_0$~: $\int_{-a}^{a} (a - |t|)/a^2 \,\mathrm{d}t =
2 \int_0^a (a - t)/a^2 \,\mathrm{d}t = 2 (a^2/2)/a^2 = 1$. La moyenne
vaut
$x_0$ par symétrie~; la variance vaut $2 \int_0^a t^2 (a - t)/a^2
\,\mathrm{d}t = \frac{2}{a^2}\left(\frac{a^4}{3} - \frac{a^4}{4}\right) =
\frac{a^2}{6}$, donc $u = a/\sqrt6 = 0.41\,a$, contre $a/\sqrt3 = 0.58\,a$
dans le cas rectangulaire~: savoir que le milieu est plus probable
réduit l'incertitude.
\end{solution}

\begin{solution}{exo:b1:lab-techniques-1:12}
$c = m/(MV) = 0.25100/(204.2 \times 0.10000) = \qty{0.012292}{mol/L}$.
\omterm{def:b1:lab-techniques-1:uncertainty}{Incertitudes relatives}~: masse $1.7 \times 10^{-4}/0.2510 = 6.8 \times
10^{-4}$, volume $0.06/100.00 = 6.0 \times 10^{-4}$~; composée~:
$\sqrt{(6.8^2 + 6.0^2)} \times 10^{-4} = 9.1 \times 10^{-4}$, soit
\qty{0.091}{\percent}~; $u(c) = \qty{1.1e-5}{mol/L}$, $U = \qty{2.2e-5}{mol/L}$~:
$c = (0.012292 \pm 0.000022)$~mol/L, $k = 2$. La pesée domine, de
peu~; les deux sources sont du même ordre.
\end{solution}

\begin{solution}{pb:b1:lab-techniques-1:1}
\textbf{1.} GHS02 inflammable, GHS05 corrosif, GHS06 toxicité aiguë,
GHS07
nocif ou irritant.
\textbf{2.} Acide salicylique~: \emph{Danger}, à cause de H318 (lésions
oculaires graves, catégorie 1). Aspirine~: \emph{Attention} (H302
seulement).
\textbf{3.} Le \omterm{def:b1:lab-techniques-1:hazard}{danger} est fixé, l'exposition ne l'est pas~: un petit
volume, mesuré et utilisé sous la hotte, avec gants, lunettes et blouse,
le flacon refermé aussitôt, aucune flamme à proximité.
\textbf{4.} H314 (graves brûlures de la peau et graves lésions des
yeux)~: lunettes et gants.
H332 (nocif par inhalation) et H226 (liquide et vapeurs inflammables)~:
la hotte et l'absence de flamme.
\textbf{5.} Rincer avec précaution à l'eau pendant plusieurs minutes, en
enlevant les lentilles de contact si possible, et continuer à rincer
(P305+P351+P338)~; consulter ensuite un médecin.
\textbf{6.} Dans les déchets aqueux, après neutralisation, et non à
l'évier.
\textbf{7.} \ce{C7H6O3 + C4H6O3 -> C9H8O4 + C2H4O2}~: C $7 + 4 = 9 + 2$,
H $6 + 6 = 8 + 4$, O $3 + 3 = 4 + 2$.
\textbf{8.} $2.00/138.0 = \qty{1.449e-2}{mol}$~; au plus
$\num{1.449e-2} \times 180.0 = \qty{2.61}{g}$ d'aspirine.
\textbf{9.} Le solide brut contient de l'acide salicylique n'ayant pas
réagi, de l'acide éthanoïque et de l'eau~: sa masse n'est pas celle de
l'aspirine, et sa pureté est inconnue.
\textbf{10.} Ce qui reste dissous après refroidissement est perdu~: une
\omterm{def:b1:precipitation:solubility}{solubilité} faible à froid et grande à chaud signifie que le produit se
dissout à chaud et revient presque entièrement à froid.
\textbf{11.} Chaque millilitre supplémentaire retient en solution, à
froid, sa part de produit.
\textbf{12.} Une croissance lente donne des \omterm{def:b1:crystals-metals:crystal}{cristaux} plus gros et plus
réguliers, qui piègent moins d'eaux mères et d'impuretés~; la glace
abaisse la \omterm{def:b1:precipitation:solubility}{solubilité}, donc la perte.
\textbf{13.} $\qty{0.045}{L} \times \qty{3.3}{g/L} = \qty{0.15}{g}$.
\textbf{14.} $1.95/2.61 = \qty{75}{\percent}$ (\qty{74.7}{\percent}).
\textbf{15.} $807/6 = \qty{134.5}{\celsius}$.
\textbf{16.} Écarts $-0.5$, $0.5$, $-1.5$, $0.5$, $-0.5$, $1.5$~; somme
des carrés 5.5~; $s = \sqrt{5.5/5} = \qty{1.05}{\celsius}$.
\textbf{17.} $u_A = 1.05/\sqrt6 = \qty{0.43}{\celsius}$.
\textbf{18.} Demi-largeur \qty{0.5}{\celsius}~: $0.5/\sqrt3 = \qty{0.29}{\celsius}$.
\textbf{19.} $1/\sqrt3 = \qty{0.58}{\celsius}$~;
$u = \sqrt{0.428^2 + 0.289^2 + 0.577^2} = \sqrt{0.600} = \qty{0.77}{\celsius}$.
\textbf{20.} $U = 2 \times 0.775 = \qty{1.5}{\celsius}$~:
$T = (134.5 \pm 1.5)$~\unit{\celsius}, $k = 2$.
\textbf{21.} Il fond bas et sur un intervalle de \qty{8}{\celsius}~: il
est impur (acide salicylique, acide éthanoïque, eau).
\textbf{22.} Donnée à l'unité, la valeur est connue à $\pm\qty{0.5}{\celsius}$ près~:
$u_{\mathrm{ref}} = 0.5/\sqrt3 = \qty{0.29}{\celsius}$.
\textbf{23.} $2.3/6.0 = 0.38$ et $3.3/6.0 = 0.55$. Le produit brut
contenait de l'acide salicylique (même $R_f$ que la référence) et de
l'aspirine~; après \omterm{def:b1:lab-techniques-1:recrystallisation}{recristallisation}, seule la tache de l'aspirine
subsiste~: l'acide salicylique a été éliminé, du moins en deçà de ce que
la plaque détecte.
\textbf{24.} L'aspirine se décompose en fondant~; chauffée lentement,
elle a le temps de se décomposer, et le mélange formé fond plus bas. La
valeur tabulée correspond à un chauffage rapide, comme sur le banc.
\textbf{25.} $z = |134.5 - 135|/\sqrt{0.600 + 0.083} = 0.5/0.827 =$
\textbf{0.60}~: bien en dessous de 2, les températures de fusion
mesurée et tabulée sont compatibles.
\end{solution}
